\documentclass[a4paper,12pt,twocolumn,landscape]{article}

\usepackage{FabZ}
\usepackage{vecteurs}
\usepackage{repere}

\usepackage{geometry}
\geometry{hmargin=0.5cm,vmargin=1.5cm}

\newcommand{\milieu}[6]{$\pointcoord{A_{\theenumi}}{#1}{#2}$ et $\pointcoord{B_{\theenumi}}{#3}{#4}$ \hfill \reponseEX{$\pointcoord{M_{\theenumi}}{#5}{#6}$}}

\newcommand{\extremite}[6]{$\pointcoord{A_{\theenumi}}{#1}{#2}$ et $\pointcoord{M_{\theenumi}}{#3}{#4}$ \hfill \reponseEX{$\pointcoord{B_{\theenumi}}{#5}{#6}$}}

\newcommand{\longueur}[5]{$\pointcoord{A_{\theenumi}}{#1}{#2}$ et $\pointcoord{B_{\theenumi}}{#3}{#4}$ \hfill \reponseEX{Longueur $A_{\theenumi}B_{\theenumi}$ = $#5$}}

\newcommand{\quadrilatere}[9]{$\pointcoord{A_{\theenumi}}{#1}{#2}$, $\pointcoord{B_\theenumi}{#3}{#4}$ $\pointcoord{C_{\theenumi}}{#5}{#6}$, $\pointcoord{D_\theenumi}{#7}{#8}$ \\ \textcolor{red}{$A_{\theenumi}B_{\theenumi}C_{\theenumi}D_{\theenumi}$ est un #9.}}

\newcommand{\parallelogrammeManqueA}[8]{$\pointcoord{B_{\theenumi}}{#1}{#2}$, $\pointcoord{C_{\theenumi}}{#3}{#4}$ $\pointcoord{D_{\theenumi}}{#5}{#6}$ \hfill \textcolor{red}{$\pointcoord{A_{\theenumi}}{#7}{#8}$}}

\newcommand{\parallelogrammeManqueB}[8]{$\pointcoord{A_{\theenumi}}{#1}{#2}$, $\pointcoord{C_{\theenumi}}{#3}{#4}$ $\pointcoord{D_{\theenumi}}{#5}{#6}$ \hfill \textcolor{red}{$\pointcoord{B_{\theenumi}}{#7}{#8}$}}

\newcommand{\parallelogrammeManqueC}[8]{$\pointcoord{A_{\theenumi}}{#1}{#2}$, $\pointcoord{B_{\theenumi}}{#3}{#4}$ $\pointcoord{D_{\theenumi}}{#5}{#6}$ \hfill \textcolor{red}{$\pointcoord{C_{\theenumi}}{#7}{#8}$}}

\newcommand{\parallelogrammeManqueD}[8]{$\pointcoord{A_{\theenumi}}{#1}{#2}$, $\pointcoord{B_{\theenumi}}{#3}{#4}$ $\pointcoord{C_{\theenumi}}{#5}{#6}$ \hfill \textcolor{red}{$\pointcoord{D_{\theenumi}}{#7}{#8}$}}

\newcommand{\rep}[3]{~$\left(#1,~#2,~#3\right)$}


%
%\setlength{\headheight}{0pt}
%\pagestyle{fancyplain}
%\fancyhf{}
%\lhead[]{\textbf{TD Vecteurs}}
%\chead[]{}
%\rhead[]{}
%
%\lfoot[]{}
%\cfoot[]{}
%\rfoot[]{}
%%\rfoot[]{Page \thepage~ sur \pageref{LastPage}}

\fancypagestyle{vide}
{
	\setlength{\headheight}{0em}
	\pagestyle{fancyplain}
	\def\headrulewidth{0em}
	\fancyhf{}
	\lhead[]{}
	\chead[]{}
	\rhead[]{}
	
	\lfoot[]{}
	\cfoot[]{}
	\rfoot[]{}
%	\rfoot[]{Page \thepage~ sur \pageref{LastPage}}
}

\usetikzlibrary{fadings}
\newcommand{\Fin}{node[xshift=-1.5ex,rotate=10]{F}
node[rotate=170]{i}
node[xshift=1.5ex,rotate=45]{n}}

\newcommand{\bonnesvacances}{node[xshift=0ex,rotate=10]{B}
node[xshift=1*1.5ex,rotate=170]{o}
node[xshift=2*1.5ex,rotate=45]{n}
node[xshift=3*1.5ex,rotate=128]{n}
node[xshift=4*1.5ex,rotate=75]{e}
node[xshift=5*1.5ex,rotate=130]{s}
node[xshift=6*1.5ex,rotate=85]{~}
node[xshift=7*1.5ex,rotate=128]{v}
node[xshift=8*1.5ex,rotate=43]{a}
node[xshift=9*1.5ex,rotate=4]{c}
node[xshift=10*1.5ex,rotate=145]{a}
node[xshift=11*1.5ex,rotate=5]{n}
node[xshift=12*1.5ex,rotate=25]{c}
node[xshift=13*1.5ex,rotate=105]{e}
node[xshift=14*1.5ex,rotate=45]{s}
}


\begin{document}


\begin{minipage}{0.45\textwidth}
\thispagestyle{firststyle}

\vspace*{1em}

\exercice
On veut que $ABCD$ soit un parallélogramme. On sait que les diagonales d'un parallélogramme ont même milieu. $[AC]$ et $[BD]$ doivent donc avoir le même milieu.\\ On peut déterminer le milieu $M$ de $[BD]$~: $\pointcoord{M}{7}{5}$.\\ Il doit être le milieu de $[AC]$ également.

Soit $\pointcoord{M}{7}{5}$ milieu de $[AC]$~:\\[1em]
\begin{minipage}{0.5\textwidth}
$
\begin{array}{rl}
	7	&=	\dfrac{ 1 + x_C}{2} \\
	14	&=	1 + x_C \\
	13	&=	x_C
  \end{array}
$
\end{minipage}
\begin{minipage}{0.5\textwidth}
$
\begin{array}{rl}
	5	&=	\dfrac{ 1 + y_C}{2} \\
	10	&=	1 + y_C \\
	9	&=	y_C
  \end{array}
$
\end{minipage}\\[1em]

Donc le point $C$ a pour coordonnées$\pointcoord{~}{13}{9}$.
\begin{center}
	\fbox{$\pointcoord{C}{13}{9}$}
\end{center}

\textbf{\underline{Remarque :} Il n'était pas demandé de placer le point $C$ ni de tracer le parallélogramme~!}

\begin{center}

\begin{tikzpicture}[scale=0.6,every node/.style={scale=0.6}]

%\coordinate (X) at (-0.5,-0.5);
%\coordinate (Y) at (13,13);
%\clip (X) rectangle (Y);
\placerpoint{A}{1}{1}{below left};
\placerpoint{B}{10}{5}{below right};
\placerpoint{D}{4}{5}{above left};
\placerpoint{C}{13}{9}{below right}
\placerpoint{M}{7}{5}{above left};
\draw[very thick,red] (A) -- (B) -- (C) -- (D) -- cycle;
\draw[very thick,red] (D) -- (M) node[midway,sloped] {$/$};
\draw[very thick,red] (B) -- (M) node[midway,sloped] {$/$};
\draw[very thick,red] (A) -- (M) node[midway,very thick,sloped] {$//$};
\draw[very thick,red] (C) -- (M) node[midway,very thick,sloped] {$//$};
\repereOIJ{-1}{15}{-1}{11};
\foreach \r in {5,10}
    \draw[thick, below right] (\r,0) node{\r};
\foreach \r in {5,10}
    \draw[thick, above left] (0,\r) node{\r};

\end{tikzpicture}

\end{center}


\end{minipage}

\vspace*{-2em}

\newpage

\begin{minipage}{0.45\textwidth}
\thispagestyle{firststyle}

\vspace*{1em}


\exercice
On se place dans le repère~\emph{orthonormé}\rep{O}{I}{J}. Dans ce repère, on considère les points $\pointcoord{A}{-3}{0}$, $\pointcoord{B}{-2}{-3}$ et $\pointcoord{C}{10}{3}$.\\Vos réponses seront argumentées.

	\begin{enumerate}
		\item $CB = \sqrt{180}$ et $CA = \sqrt{178}$.\\
				$CB \neq CA$ donc le point $A$ n'appartient pas au cercle $\mathscr{C}$ de centre $C$ passant par~$B$.
		\item De même, $CB \neq CA$ donc le point $C$ n'appartient pas à la médiatrice de $\left[AB\right]$.
%		\item Déterminer les coordonnées de $Z$ pour que $AZBC$ soit un parallélogramme.
	\end{enumerate}


\exercice
\begin{enumerate}
	\item Dans le repère\rep{O}{I}{J}~orthonormé, on calcule~:\\ $AB=\sqrt{\left( - 28 - 122  \right) ^ 2 + \left(-98 - \left( -48 \right) \right) ^ 2}$\\
			$AB=\sqrt{\left( - 28 - 122 \right) ^ 2 + \left(-98 + 48 \right) ^ 2}$\\
			$AB=\sqrt{\left( - 150 \right) ^ 2 + \left( - 50 \right) ^ 2}$\\
			$AB=\sqrt{22500 + 2500}$\\
			$AB=\sqrt{25000}$\\[1em]
			De même, on obtient~:\\ $AC = \sqrt{25000}$ et $BC = \sqrt{50000}$\\[1em]
			On a $AB$ = $AC$ donc le triangle $ABC$ est \textbf{isocèle}.\\[1em]
			De plus, $BC^2 = \left(\sqrt{50000}\right)^2 = 50000$ et\\ $AB^2 + AC^2 = \left(\sqrt{25000}\right)^2 + \left(\sqrt{25000}\right)^2 = 50000$\\
			D'où, $BC^2 = AB^2 + AC^2$\\
			D'après la réciproque du théorème de Pythagore, le triangle $ABC$ est \textbf{rectangle} en $A$.\\[1em]
			Conclusion~: \fbox{Le triangle $ABC$ est \textbf{rectangle isocèle en $A$}.}
\end{enumerate}

\end{minipage}

\newpage

\begin{minipage}{0.45\textwidth}
\thispagestyle{vide}

\vspace*{1em}


\begin{enumerate}
	\setcounter{enumi}{1}
	\item Dans le repère\rep{O}{I}{J},\\[1em]
			$x_D = \dfrac{122 + (-28)}{2} = \dfrac{94}{2} = 47$ et\\
			$y_D = \dfrac{-98 + (-48)}{2} = \dfrac{-146}{2} = -73$\\[1em]
			Donc \fbox{$\pointcoord{D}{47}{-73}$}\\[1em]
			De même, on a~:\\[1em] \fbox{$\pointcoord{E}{22}{2}$} et \fbox{$\pointcoord{F}{97}{27}$}
	\item Un carré est un parallélogramme particulier, commençons par montrer que $ADEF$ est un parallélogramme~:\\
			Soit $M$ milieu de $[AE]$, $\pointcoord{M}{72}{-23}$\\
			Soit $N$ milieu de $[DF]$, $\pointcoord{N}{72}{-23}$\\
			$M$ et $N$ sont \textbf{confondus}.\\
			Un quadrilatère dont les diagonales ont leurs milieux \textbf{confondus} est un parallélogramme. $M$ et $N$ sont confondus donc le quadrilatère $ADEF$ est un parallélogramme.\\[1em]
			$AE = \sqrt{12500}$ et $DF = \sqrt{12500}$\\
			Un parallélogramme dont les \textbf{diagonales ont même longueur} est un rectangle. $AE = DF$ donc le parallélogramme $ADEF$ est un rectangle.\\[1em]
			$AD = \sqrt{6250}$ et $DE = \sqrt{6250}$\\			
			Un rectangle qui a \textbf{deux côtés consécutifs de même longueur} est un carré. $AD = DE$ donc le rectangle $ADEF$ est un carré.\\
			\fbox{Donc le quadrilatère $ADEF$ est un carré.}\\
	\item D'après la question précédente, on sait que $ADEF$ est un carré donc en particulier que $DE = EF$ et que $\widehat{DEF} = 90°$.\\[1em]
			\fbox{Donc le triangle $DEF$ est \textbf{rectangle isocèle}.}
\end{enumerate}

\end{minipage}

\newpage

\begin{minipage}{0.45\textwidth}
\thispagestyle{firststyle}

\vspace*{1em}


\begin{enumerate}
	\setcounter{enumi}{4}
	\item Dans le \emph{repère\rep{A}{B}{C}}, $\pointcoord{A}{0}{0}$ ; $\pointcoord{B}{1}{0}$ ; $\pointcoord{C}{0}{1}$ ; $\pointcoord{D}{\dfrac{1}{2}}{0}$ ; $\pointcoord{E}{\dfrac{1}{2}}{\dfrac{1}{2}}$ et $\pointcoord{F}{0}{\dfrac{1}{2}}$.
	\item Le repère\rep{A}{B}{C} est \textbf{orthonormé} ($AB = AC$ et $\widehat{BAC} = 90°$ OU $ABC$ rectangle isocèle en $A$) \textbf{donc} \textbf{ON PEUT} déterminer la longueur $AB$ dans ce repère.
\end{enumerate}


\vspace*{1em}
\centering
\begin{tikzpicture}[scale=2,transform shape]
\draw (0,0) \Fin;
\draw (-1em,-1ex) -- (1em,-1ex);
\path[scope fading=south] (-1em,-0.25em) rectangle (1em,-3.75ex);
\draw[yscale=-1] (0,2ex) \Fin;
\end{tikzpicture}
\vspace*{\stretch{1}}

\end{minipage}

\end{document}
